\documentclass[11pt,letterpaper]{article}
\usepackage[margin=1in]{geometry}
\usepackage[T1]{fontenc}
\usepackage{graphicx,booktabs,amsmath}
\usepackage{xcolor}
\usepackage{hyperref}
\usepackage{natbib}
\usepackage{titlesec}
\usepackage{fancyhdr}
\hypersetup{colorlinks=true,linkcolor=blue!70!black,urlcolor=blue!70!black,citecolor=blue!70!black,pdftitle={ECE750T Assignment 2: Florence-2 Documentation Audit},pdfauthor={Instructor example}}
\titleformat{\section}{\normalfont\normalsize\bfseries}{}{0pt}{}
\titlespacing*{\section}{0pt}{10pt}{4pt}
\setlength{\parindent}{0pt}
\setlength{\parskip}{6pt}
\pagestyle{fancy}
\fancyhf{}
\renewcommand{\headrulewidth}{0pt}
\fancyfoot[R]{\footnotesize\thepage}

\begin{document}
{\Large\textbf{ECE750T: Data Documentation Audit}}\\[5pt]
{\large\textbf{Florence-2-base}}\\[4pt]
Example submission \quad | \quad Student ID: not applicable\\
Fall 2026
\vspace{0.5em}\hrule\vspace{0.5em}

\section*{1. Context [8 marks]}
I audit Microsoft's Hugging Face model card for \texttt{microsoft/Florence-2-base}, a 230-million-parameter vision model. The artifact is the repository README at \url{https://huggingface.co/microsoft/Florence-2-base}, accessed October 13, 2026, revision \texttt{5ca5edf5} (full revision in the references). I use its linked technical report as supporting evidence, distinguishing information in that report from information in the card \citep{microsoftCard,xiao2023}.

The proposed use is drafting English descriptions of informative campus photographs for a university website; editors review drafts in their page context before publication for screen-reader users. This is my assessment scenario, not an established deployment. The card presents a general-purpose model for visual tasks, including captioning, detection and OCR, rather than an accessibility-specific tool. Its examples and benchmarks suggest an audience of developers and researchers. Its purpose is to explain capabilities and enable use; this audit asks whether it also enables a website team to judge suitability. I assess it using the Model Cards framework \citep{mitchell2019}.

\section*{2. Critique [20 marks]}
\textbf{Strengths.} The card's ``Model Summary'' distinguishes base, large and fine-tuned variants; ``Tasks'' specifies prompts and output formats; ``Benchmarks'' names datasets and metrics. These details support identifying the correct artifact and trying a relevant task. For example, the base model's reported COCO caption score is 133.0 CIDEr under ``Florence-2 Zero-shot performance.'' Scores for \texttt{base-ft} should not be attributed to the base checkpoint. Linking the paper also gives readers a route to inspect methodology.

\textbf{Use boundaries are unclear.} The card does not explicitly identify primary users, out-of-scope uses or deployment caveats, as recommended by \citet[\S\S4.2, 4.8--4.9]{mitchell2019}. A list of supported tasks does not explain whether drafts require checking for invented details, misleading descriptions of people or incorrect visible text. These are plausible error pathways, not measured findings from this audit. Without guidance, an editor cannot judge which outputs to reject or when to write a description manually.

\textbf{The evaluation does not establish accessibility usefulness.} CIDEr compares generated captions with human reference descriptions \citep{vedantam2015}. That is useful evidence of captioning performance, but it does not directly measure whether a reader understands the photograph's purpose on a particular page. W3C guidance makes that purpose central and permits empty alternatives for decorative or redundant images \citep{w3c2024}. Thus a fluent, reference-like caption can still omit necessary context or add irrelevant detail. Treating benchmark consensus as sufficient evidence of accessible communication would leave a construct-validity gap \citep[\S3.2]{jacobs2021}. The card supplies no study of this proposed editorial workflow or evidence that human review catches its errors.

\newpage
\textbf{Coverage and label quality remain difficult to judge.} Aggregate benchmark scores do not show errors across lighting, image quality, assistive devices or relevant populations, although the Model Cards framework recommends reporting salient factors and disaggregated results \citep[\S\S4.3, 4.7]{mitchell2019}. The paper does identify source datasets and automated annotation, filtering and refinement (\S\S4.1--4.3); it would be incorrect to claim that training provenance is entirely absent. However, neither the card nor those descriptions establishes error rates for descriptions of our campus population. Automated labels could reproduce an annotator's omissions; the documentation does not resolve this possibility. Missing evidence establishes uncertainty, not demonstrated discrimination. The card also omits uncertainty estimates and a complete benchmark reproduction recipe, limiting independent checks on its aggregate claims.

\section*{3. Proposed changes [12 marks]}
\textbf{First, add a use-and-limitations section.} Explain that published benchmarks do not establish suitability for this accessibility workflow. Recommend checking every draft against the image and page purpose, including visible text and unsupported personal inferences; describe when to replace it manually. Microsoft maintainers should state evaluated uses and known limitations. Website editors and accessibility staff should define review responsibilities and publication criteria. These additions would turn broad task claims into actionable guidance without inventing evidence.

\textbf{Second, publish an evaluation relevant to the decision.} With consenting screen-reader users and accessibility specialists, define faithful and useful descriptions, then evaluate representative campus photographs in page context. Compare model-assisted editing with manual writing; report unsupported-claim and critical-omission rates, user task success, and editing effort. Include challenging capture conditions and ethically justified group comparisons, without inferring sensitive identities from photographs. Report sampling, label definitions, evaluator agreement, subgroup counts and uncertainty; mark small or untested groups as unresolved. The website team can supply permitted images and page purposes; users and specialists can supply judgments; Microsoft or independent evaluators can run the assessment. Do not claim validation until those results exist.

\textbf{Third, make evidence traceable.} Add a concise training-data coverage and annotation-quality summary, linked to the paper, stating what remains unknown. Microsoft's data team could supply source proportions, filtering decisions and human validation of automated labels. For each evaluation, maintainers should record checkpoint and documentation revisions, dataset splits, preprocessing, prompts, decoding settings and metric implementation, with a changelog and feedback contact. This would let readers reproduce relevant claims and track whether limitations change. Until the missing evidence is supplied, the card supports a supervised pilot, not a conclusion that the proposed workflow is suitable.

\clearpage
\begingroup
\raggedright
\bibliographystyle{plainnat}
\bibliography{references}
\endgroup

\section*{AI-assistance disclosure and example provenance}
This is an instructor-requested, AI-assisted example for distribution after the assignment, formatted as a student submission. Codex assisted with source research, analysis, drafting, LaTeX preparation and submission-file generation. No model inference or user study was conducted. The identifying header denotes an example rather than a real student. The assignment permits students to use AI only to edit their own prose, with disclosure; this instructor example does not change that policy.
\end{document}
